本节给出模板中常用排版方式的简要说明，供写作时参考。

\subsection{公式与编号}
这一部分可以按“规范 - 示例 - 常见错误”的顺序来讲。
常用的公式环境有 \texttt{equation}、\texttt{align} 和 \texttt{subequations}。
单行公式用 \texttt{equation}，多行推导用 \texttt{align}，相关公式组可用 \texttt{subequations} 统一编号。

\begin{equation}\label{eq:simple}
E = mc^2
\end{equation}

\begin{align}
a &= b + c \label{eq:align-a}\\
  &= d + e. \label{eq:align-b}
\end{align}

写公式时建议遵循“先编号，后引用”的顺序。
\texttt{\textbackslash ref}、\texttt{\textbackslash eqref} 和 \texttt{\textbackslash cref} 的区别如下：
\begin{itemize}
    \item \texttt{\textbackslash ref\{eq:simple\}} 只输出编号本身，例如“1”。
    \item \texttt{\textbackslash eqref\{eq:simple\}} 会自动加上圆括号，例如“(1)”。
    \item \texttt{\textbackslash cref\{eq:simple\}} 会自动带上对象名称，例如\cref{eq:simple}。如果一次引用多个标签，也会自动处理成更自然的连写形式。
\end{itemize}
一般优先使用 \texttt{\textbackslash cref}。

\paragraph{align 与 equation+aligned 的区别}
\texttt{align} 适合多行推导，并且每一行都可以单独编号；如果某一行不需要编号，可以用 \texttt{\textbackslash nonumber} 或 \texttt{\textbackslash notag} 取消。
这种写法适合中间步骤也需要单独引用的场景。

\begin{align}
f(x) &= x^2 + 1 \label{eq:align-demo-a}\\
     &= (x+1)^2 - 2x. \label{eq:align-demo-b}
\end{align}

如果只是想把一个公式块排成多行，但仍然只给整个式子一个编号，可以在 \texttt{equation} 里嵌套 \texttt{aligned}。
这种写法只产生一个编号，更适合“一个结论、多个对齐步骤”的情况。

\begin{equation}\label{eq:equation-aligned}
\begin{aligned}
f(x) &= x^2 + 1\\
     &= (x+1)^2 - 2x.
\end{aligned}
\end{equation}

简要地说：
\begin{itemize}
    \item \texttt{align} = 多行推导，每行可独立编号。
    \item \texttt{equation + aligned} = 一个编号的多行排版。
    \item 如果只是显示公式、不需要编号，可用 \texttt{\textbackslash[ ... \textbackslash ]}。
\end{itemize}

\[
f(x) = x^2 + 1
\]

\texttt{\textbackslash[ ... \textbackslash ]} 与 \texttt{equation} 的视觉效果接近，都是独立居中的显示公式；区别是前者默认不编号，后者默认编号。
因此，若公式不需要引用，优先用 \texttt{\textbackslash[ ... \textbackslash ]}；若需要编号再引用，则用 \texttt{equation} 或 \texttt{align}。

\paragraph{常见错误}
\begin{enumerate}
    \item \textbf{把普通文本直接放进数学模式。} 数学模式只适合变量、算子和公式结构，不适合写完整句子。

    \[
    \text{This is text}
    \]

    \item \textbf{把中文句子直接写进公式。} 公式环境中尽量只保留数学内容。

    \[
    \text{这是一个公式}
    \]

    \item \textbf{多行公式不用 \texttt{align}。} 需要对齐时，应使用 \texttt{align}，不要把每一行直接分开写。

    \begin{align}
    a &= b + c \\
    d &= e + f
    \end{align}

    \item \textbf{微分写法不规范。} 建议写成 \texttt{\textbackslash int f(x)\textbackslash, \textbackslash mathrm\{d\}x}，不要直接写成 \texttt{dx}。

    \[
    \int f(x)\,\mathrm{d}x
    \]
    \item \textbf{正体使用不规范。} 变量用斜体，函数名、单位和常数用正体。比如$\mathrm{e}$、$\sin$、$\mathrm{kg}$等。
    \item \textbf{ \texttt{\$\$...\$\$}是老版和Markdown用法。} 正式论文中不推荐使用 \texttt{\$\$}，应改用 \texttt{\textbackslash[...\textbackslash]} 或 \texttt{equation}。
    \item \textbf{空格不会用。} 数学模式中的空格需要用 \texttt{\textbackslash,}、\texttt{\textbackslash;}、\texttt{\textbackslash quad} 等命令来控制，直接敲空格键是无效的。

    \item \textbf{括号大小不匹配。} 分式较长时应使用自动伸缩括号。

    \[
    \left( \frac{a}{b} \right)
    \]

    \item \textbf{手工编号。} 不要直接写成 \texttt{(1)}，应使用 \texttt{\textbackslash label} 和 \texttt{\textbackslash ref} 自动引用。

    \item \textbf{忽略特殊字符转义。} 例如 \texttt{\%}、\texttt{\$}、\texttt{\_}、\texttt{\&} 都需要转义后才能正常输出。

\end{enumerate}

\subsection{表格排版}
表格建议使用 \texttt{booktabs} 的 \texttt{\textbackslash toprule}、\texttt{\textbackslash midrule}、\texttt{\textbackslash bottomrule}。
尽量少用竖线，表头放在表格上方。

\begin{table}[H]
    \centering
    \caption{表格排版示意}
    \label{tab:layout}
    \begin{tabular}{cc}
        \toprule
        符号 & 含义 \\
        \midrule
        $x$ & 变量 \\
        $y$ & 观测值 \\
        \bottomrule
    \end{tabular}
\end{table}

正文中可写成表\cref{tab:layout}。

\subsection{算法排版}
如果需要展示求解步骤，建议使用 \texttt{algorithm} 和 \texttt{algorithmic}。
这比直接堆列表更清晰，也更适合“输入 - 处理 - 输出”的结构。

\begin{algorithm}[H]
\caption{示意算法}
\label{alg:demo}
\begin{algorithmic}
\STATE 输入：初始值 $x_0$
\STATE 输出：结果
\STATE 计算目标函数
\WHILE{未收敛}
    \STATE 更新变量
\ENDWHILE

\end{algorithmic}
\end{algorithm}

正文中可写成算法\cref{alg:demo}。

\subsection{定理与证明}
模板里已经定义了定理、引理、推论和证明环境。
适合把结论、约束和关键推导分开写，逻辑会更清楚。

\begin{theorem}
若 $A$ 为对称矩阵，则其特征值均为实数。
\end{theorem}

\begin{proof}
由对称矩阵的标准性质可知结论成立。
\end{proof}

\subsection{结构与附录}
正文通常按“问题重述、问题分析、模型建立、问题求解、模型评价”的顺序组织。
附录适合放推导细节、补充计算和完整代码，主文中只保留必要结果。

\subsection{参考文献与交叉引用}
正文中使用 \texttt{\textbackslash cite} 引用文献，文末由 \texttt{biblatex}/\texttt{biber} 根据 \texttt{reference.bib} 自动生成参考文献\cite{nan2017fast}。
图、表、公式、算法最好都统一使用 \texttt{\textbackslash label} 配合 \texttt{\textbackslash cref} 或 \texttt{\textbackslash ref} 引用，避免手工编号出错。
